% Section 2: Distribution Learning & Probability Foundations

\section{Distribution Learning \& Probability}

\begin{frame}[fragile]{Empirical Data Distribution vs Parametric Density}
\begin{columns}
\column{0.48\textwidth}
\textbf{Mathematical Formulation}

Given observed image samples $\{x_1, x_2, \dots, x_N\} \sim p_{\text{data}}(x)$:
\begin{equation*}
p_{\text{data}}(x) = \frac{1}{N} \sum_{i=1}^N \delta(x - x_i)
\end{equation*}

Our goal is to learn a continuous parametric distribution $p_\theta(x)$ that approximates $p_{\text{data}}(x)$.

Generative modeling = matching distributions:
\begin{equation*}
p_\theta(x) \approx p_{\text{data}}(x)
\end{equation*}

\vspace{1mm}
\textbf{\color{commentgreen}Biological Intuition (DAPI $\to$ GFP):}\\
\fontsize{6.5pt}{7.5pt}\selectfont
$p_{\text{data}}$ represents the real cell population containing a mixture of GFP-positive (bright) and GFP-negative (dark) cells. The model distribution $p_\theta$ must match this biological mixture.

\column{0.52\textwidth}
\textbf{Python Code Equivalent}
\begin{lstlisting}
import torch
import torch.distributions as dist

# Empirical sample batch x_data
x_data = torch.randn(100, 1) * 2.0 + 5.0

# Parametric model distribution p_theta = N(mu, sigma)
mu = torch.tensor([0.0], requires_grad=True)
sigma = torch.tensor([1.0], requires_grad=True)

# Define parametric distribution family
p_theta = dist.Normal(loc=mu, scale=sigma)
\end{lstlisting}
\end{columns}
\end{frame}


\begin{frame}[fragile]{Maximum Likelihood Estimation (MLE)}
\begin{columns}
\column{0.48\textwidth}
\textbf{Mathematical Formulation}

Finding optimal parameters $\theta^*$ by maximizing log-likelihood:
\begin{equation*}
\theta^* = \arg\max_\theta \sum_{i=1}^N \log p_\theta(x_i)
\end{equation*}

Equivalent to minimizing Negative Log-Likelihood (NLL) / Forward KL:
\begin{equation*}
\mathcal{L}_{\text{NLL}}(\theta) = -\E_{x \sim p_{\text{data}}} \big[ \log p_\theta(x) \big]
\end{equation*}

\vspace{1mm}
\textbf{\color{commentgreen}Biological Intuition (DAPI $\to$ GFP):}\\
\fontsize{6.5pt}{7.5pt}\selectfont
Optimizes network parameters so that generating bright GFP signals for positive cells (and zero signal for negative cells) has maximum probability under true micrographs.

\column{0.52\textwidth}
\textbf{Python Code Equivalent}
\begin{lstlisting}
def fit_mle(x_samples, p_theta_cls, lr=0.1, steps=100):
    mu = torch.tensor([0.0], requires_grad=True)
    sigma = torch.tensor([1.0], requires_grad=True)
    opt = torch.optim.Adam([mu, sigma], lr=lr)
    
    for _ in range(steps):
        opt.zero_grad()
        p_theta = p_theta_cls(loc=mu, scale=F.softplus(sigma))
        
        # Negative Log-Likelihood Loss
        nll_loss = -p_theta.log_prob(x_samples).mean()
        nll_loss.backward()
        opt.step()
        
    return mu, F.softplus(sigma)
\end{lstlisting}
\end{columns}
\end{frame}


\begin{frame}[fragile]{Kullback-Leibler (KL) Divergence}
\begin{columns}
\column{0.48\textwidth}
\textbf{Mathematical Formulation}

Asymmetry measure of distance between two probability densities $p$ and $q$:
\begin{equation*}
\KL{p}{q} = \int p(x) \log \frac{p(x)}{q(x)} \, dx
\end{equation*}

Analytical solution for two Gaussians $\mathcal{N}_0(\mu_0, \sigma_0^2)$ and $\mathcal{N}_1(\mu_1, \sigma_1^2)$:
\begin{equation*}
\KL{p_0}{p_1} = \log \frac{\sigma_1}{\sigma_0} + \frac{\sigma_0^2 + (\mu_0 - \mu_1)^2}{2\sigma_1^2} - \frac{1}{2}
\end{equation*}

\vspace{1mm}
\textbf{\color{commentgreen}Biological Intuition (DAPI $\to$ GFP):}\\
\fontsize{6.5pt}{7.5pt}\selectfont
Quantifies information loss when approximating true GFP brightness distributions. Forward KL heavily penalizes predicting false GFP expression in cells that are actually negative.

\column{0.52\textwidth}
\textbf{Python Code Equivalent}
\begin{lstlisting}
# Analytical KL Divergence between 2 Gaussians in PyTorch

def kl_divergence_gaussians(mu0, log_var0, mu1=0.0, log_var1=0.0):
    # KL( N(mu0, var0) || N(mu1, var1) )
    var0 = torch.exp(log_var0)
    var1 = torch.exp(torch.tensor(log_var1))
    
    kl = 0.5 * (
        (log_var1 - log_var0) + 
        (var0 + (mu0 - mu1)**2) / var1 - 
        1.0
    )
    return kl.sum(dim=-1)
\end{lstlisting}
\end{columns}
\end{frame}


\begin{frame}[fragile]{Wasserstein Distance (Optimal Transport)}
\begin{columns}
\column{0.48\textwidth}
\textbf{Mathematical Formulation}

Earth Mover's Distance ($W_1$ metric):
\begin{equation*}
W_1(p, q) = \inf_{\gamma \in \Pi(p,q)} \E_{(x,y)\sim\gamma} \big[ \|x - y\| \big]
\end{equation*}

Kantorovich-Rubinstein Duality:
\begin{equation*}
W_1(p_{\text{data}}, p_\theta) = \sup_{\|f\|_L \le 1} \E_{x \sim p_{\text{data}}}[f(x)] - \E_{x \sim p_\theta}[f(x)]
\end{equation*}

Key benefit: Smooth gradients even when support of $p$ and $q$ do not overlap!

\vspace{1mm}
\textbf{\color{commentgreen}Biological Intuition (DAPI $\to$ GFP):}\\
\fontsize{6.5pt}{7.5pt}\selectfont
Measures the minimum "work" needed to shift predicted cell GFP fluorescence profiles to match true biological GFP distributions across positive and negative populations.

\column{0.52\textwidth}
\textbf{Python Code Equivalent}
\begin{lstlisting}
# 1D Empirical Wasserstein-1 Distance calculation

def compute_w1_distance(x_real, x_gen):
    # Sort samples along batch dimension
    x_real_sorted, _ = torch.sort(x_real, dim=0)
    x_gen_sorted, _ = torch.sort(x_gen, dim=0)
    
    # EMD in 1D is L1 distance between sorted quantiles
    w1_dist = torch.mean(torch.abs(x_real_sorted - x_gen_sorted))
    return w1_dist
\end{lstlisting}
\end{columns}
\end{frame}
